\begin{table}[!htbp]
\caption{DELB 四部分定理体系及各部分回答的不同问题。贡献在于针对因果整数
预测的联合组织；基础最大熵和信息投影结果仍属于既有理论。}
\label{tab:delb_theorem_package}
\centering
\footnotesize
\setlength{\tabcolsep}{4pt}
\begin{tabularx}{\textwidth}{
    >{\RaggedRight\arraybackslash}p{0.19\textwidth}
    >{\RaggedRight\arraybackslash}p{0.28\textwidth}
    >{\RaggedRight\arraybackslash}X}
\toprule
\textbf{组成} & \textbf{数学对象} & \textbf{回答的预测问题}\\
\midrule
精确逆格点底线
&
\(\mathcal H_{\mathbb Z}^{-1}
 [H(X_t\mid\mathcal F_{t-m})]\)
&
目标不确定性和整数格点对所有算法施加何种 MSE 下界？\\

预测器相关强化
&
\(\mathcal H_{\mathbb Z}^{-1}
 [H(X_t\mid\mathcal F_{t-m})+I(E_t;\mathcal F_{t-m})]\)
&
特定预测器的误差仍保留历史信息时，下界如何加强？\\

可达充要条件
&
\(P_E=q_D\) 且 \(I(E;\mathcal F)=0\)
&
预测器何时达到普适 DELB，而不仅仅是满足该不等式？\\

双松弛诊断
&
\(I(E;\mathcal F)+D_{\mathrm{KL},2}(P_E\Vert q_D)\)
&
不可达来自残余历史依赖、误差形状失配，还是二者兼有？\\
\bottomrule
\end{tabularx}
\end{table}
